% Formato PGC — LaTeX conversion (1:1 from Formato PGC.docx)
% Build: latexmk -pdf main.tex  or  pdflatex main.tex
\documentclass[11pt]{article}

\usepackage[letterpaper,top=1.15cm,bottom=2.79cm,left=2.82cm,right=2.79cm,headheight=2.55cm,headsep=0.35cm,footskip=1.45cm,includehead,includefoot]{geometry}
\usepackage[spanish]{babel}
\usepackage[utf8]{inputenc}
\usepackage[T1]{fontenc}
\usepackage{helvet}
\renewcommand{\familydefault}{\sfdefault}
\usepackage{graphicx}
\usepackage{xcolor}
\usepackage{array}
\usepackage{colortbl}
\usepackage{multirow}
\usepackage{fancyhdr}
\usepackage{lastpage}
\usepackage{eso-pic}
\usepackage{tikz}
\usepackage{amssymb}
\usepackage{hyperref}
\usepackage{booktabs}

\definecolor{udecgreen}{HTML}{6B8E00}
\definecolor{udecgreenlight}{HTML}{8BA817}
\definecolor{udecgold}{HTML}{C5A000}

\pagestyle{fancy}
\fancyhf{}
\renewcommand{\headrulewidth}{0pt}
\renewcommand{\footrulewidth}{0pt}

% Footer: centered "1" + bottom green bar + triangle
% We use eso-pic for the green bar/triangle so it matches Word's behindDoc anchors
% and fancyfoot for the centered page number in a white box.
\fancyfoot[C]{\raisebox{0.35cm}{\colorbox{white}{\makebox[0.95cm][c]{\small\color{black}\thepage}}}}
% Header: UDeC banner + thin line — only on first page (Word sectPr 1)
\fancypagestyle{firstpage}{%
  \fancyhf{}%
  \fancyhead[C]{%
    \begin{tikzpicture}[remember picture,overlay]
      % banner image at top (behindDoc equivalent)
      \node[anchor=north west,inner sep=0] at ([xshift=0.70cm,yshift=-0.34cm]current page.north west)
        {\includegraphics[width=7.35cm,height=2.37cm,keepaspectratio]{assets/logo-universidad.jpg}};
      % thin horizontal line below banner (image2 ~ 7.7cm wide, ~0.7cm tall) — replicate as rule
      % position: just below header block, spanning content width
      \draw[udecgreen,line width=0.45pt] ([xshift=0.70cm,yshift=-2.85cm]current page.north west) -- ++(18.05cm,0);
    \end{tikzpicture}%
    \vspace{2.05cm}%
  }%
  \fancyfoot[C]{\raisebox{0.35cm}{\colorbox{white}{\makebox[0.95cm][c]{\small\color{black}\thepage}}}}%
  \renewcommand{\headrulewidth}{0pt}%
  \renewcommand{\footrulewidth}{0pt}%
}

% Green footer bar (Word footer1 anchor: behindDoc image3 + page-wide rectangle)
\AddToShipoutPictureBG{%
  \AtPageLowerLeft{%
    \begin{tikzpicture}[overlay]
      % full-width green gradient bar at bottom (~1.15cm tall)
      \fill[udecgreen] (0,0) rectangle (\paperwidth,1.18cm);
      % diagonal triangle (image3) at bottom-right — replicate with polygon fill
      \fill[white] (\paperwidth,0) -- (\paperwidth,1.18cm) -- ({\paperwidth-4.0cm},0) -- cycle;
      % inner green triangle cut
      \fill[udecgreen] ({\paperwidth-3.85cm},0) -- (\paperwidth,0) -- (\paperwidth,1.08cm) -- cycle;
    \end{tikzpicture}%
  }%
}

\setlength{\arrayrulewidth}{0.5pt}
\renewcommand{\arraystretch}{1.08}
\setlength{\tabcolsep}{5pt}
\newcolumntype{Y}{>{\raggedright\arraybackslash}X}
\newcolumntype{C}{>{\centering\arraybackslash}X}

\newcommand{\formlabel}[1]{\textbf{\small #1}}
\newcommand{\formcell}[1]{\small #1}

\begin{document}
\thispagestyle{firstpage}

% Title block: matches Word's "Anexo 1. Formato de Propuesta" (indented left) + centered "CICLO"
\vspace{0.10cm}
{\small\bfseries\hspace{0.62cm}Anexo\hspace{0.15em}1.\hspace{0.25em}Formato\hspace{0.15em}de\hspace{0.15em}Propuesta\par}
\vspace{0.20cm}
{\centering\small\bfseries CICLO II — PROFUNDIZACIÓN\par}
\vspace{0.18cm}
{\small\bfseries\hspace{0.62cm}DATOS\hspace{0.15em}GENERALES\par}
\vspace{0.12cm}

% --- DATOS GENERALES (Table 0) — two-column gridCols 4131+4131 ---
\noindent
\begin{tabular}{|p{7.55cm}|p{7.55cm}|}
\hline
\formlabel{PROGRAMA} & \formcell{Ingeniería de Sistemas y Computación} \\
\hline
\formlabel{SEDE} & \formcell{Facatativá} \\
\hline
\formlabel{TIPO DE PROYECTO} & \formcell{Proyecto Tecnológico} \\
\hline
\formlabel{CICLO} & \formcell{Ciclo II — Profundización} \\
\hline
\multicolumn{2}{|p{15.38cm}|}{\formlabel{INTEGRANTES DEL PROYECTO}} \\
\hline
\formlabel{NOMBRES Y APELLIDOS} & \formlabel{CORREO ELECTRÓNICO INSTITUCIONAL} \\
\hline
\formcell{David Fernando Gonzalez Soler} & \formcell{davidfernandogonzalez@ucundinamarca.edu.co} \\
\hline
\formcell{Maicol Estiven Fonseca Muñoz} & \formcell{mefonseca@ucundinamarca.edu.co} \\
\hline
\formcell{Ivan David Giraldo Garavito} & \formcell{idgiraldo@ucundinamarca.edu.co} \\
\hline
\end{tabular}

\vspace{0.22cm}
{\small\bfseries\hspace{0.62cm}ESTRUCTURA\hspace{0.20em}DEL\hspace{0.20em}PROYECTO\par}
\vspace{0.14cm}

% --- ESTRUCTURA (Table 1) — single column 8262 ---
\noindent
\begin{tabular}{|p{15.38cm}|}
\hline
\formlabel{TÍTULO DEL PROYECTO:} \\
\hline
\rule{0pt}{0.62cm} \\
\hline
\formlabel{PLANTEAMIENTO DEL PROBLEMA Y JUSTIFICACIÓN:} \\
\hline
\formcell{El control de los gastos personales sigue siendo una tarea manual: registrar cada compra en una libreta, una hoja de cálculo o una aplicación exige disciplina y tiempo, y la mayoría abandona el registro a las pocas semanas.

En Colombia, la vida financiera ocurre cada vez más en el celular: el 66,6 \% de las operaciones monetarias se realiza por canales digitales (Banca de las Oportunidades, 2024). Aun así, las aplicaciones de finanzas personales fracasan: el 62 \% de los usuarios las abandona durante el primer mes (Apptopia, 2024, citado en Elsalaymeh, 2026), mientras que formar un hábito requiere en promedio entre 59 y 66 días (Singh et al., 2024).

Sin registro sistemático, el usuario no sabe cuánto gasta en transporte, comida o suscripciones, no compara su gasto real contra su presupuesto y termina decidiendo a ciegas: sobregiros, crédito costoso o imposibilidad de ahorrar.

¿Cómo puede una aplicación móvil con asistencia de IA facilitar el registro y el análisis de los gastos personales, reduciendo el esfuerzo de captura y ayudando a formar el hábito de control financiero?} \\
\hline
\formlabel{OBJETIVOS DEL PROYECTO} \\
\hline
\formlabel{OBJETIVO GENERAL:} \\
\hline
\rule{0pt}{0.62cm} \\
\hline
\formlabel{OBJETIVOS ESPECÍFICOS:} \\
\hline
\rule{0pt}{0.62cm} \\
\hline
\formlabel{MARCO REFERENCIAL (Opcional):} \\
\hline
\rule{0pt}{0.62cm} \\
\hline
\formlabel{METODOLOGÍA A UTILIZAR Y CRONOGRAMA DE DESARROLLO:} \\
\hline
\rule{0pt}{0.62cm} \\
\hline
\formlabel{REFERENCIAS BIBLIOGRÁFICAS (opcional):} \\
\hline
\rule{0pt}{0.62cm} \\
\hline
\end{tabular}

\vspace{0.20cm}
% --- PRODUCTOS ESPERADOS (Table 2) — cols 562 + 7701 ---
\noindent
\begin{tabular}{|p{0.72cm}|p{14.50cm}|}
\hline
\multicolumn{2}{|p{15.38cm}|}{\formlabel{PRODUCTOS ESPERADOS (marque con una X):}} \\
\hline
$\square$ & \formcell{Documento síntesis de PGC} \\
\hline
$\square$ & \formcell{Software} \\
\hline
$\square$ & \formcell{Sistema de información} \\
\hline
$\square$ & \formcell{Prototipo} \\
\hline
$\square$ & \formcell{Artículo científico} \\
\hline
$\square$ & \formcell{Póster- Ponencia} \\
\hline
$\square$ & \formcell{Comunidad de práctica y redes de conocimiento} \\
\hline
$\square$ & \formcell{Evento científico o académico} \\
\hline
$\square$ & \formcell{Ediciones de revista o libro de divulgación} \\
\hline
$\square$ & \formcell{Producción de contenidos digitales} \\
\hline
$\square$ & \formcell{Publicaciones editoriales no especializadas} \\
\hline
$\square$ & \formcell{Desarrollo Web} \\
\hline
$\square$ & \formcell{Creación de cursos} \\
\hline
$\square$ & \formcell{Emprendimiento e innovación} \\
\hline
$\square$ & \formcell{Otro ¿Cuál?} \\
\hline
\end{tabular}

\end{document}
